\section{Motivation}
\label{sec:motivation}

We first examine what existing Korean benchmarks demand of a model and where models fail on them.
We analyze 7,340 multiple-choice questions from CLIcK~\citep{kim2024click}, KMMLU-Redux, and KMMLU-Pro~\citep{hong-etal-2025-kmmlu}.

\subsection{Silver Reasoning}
\label{sec:silver}

Identifying the knowledge a question requires and how that knowledge must be applied is challenging for an LLM annotator given only the question, since LLMs tend to underperform in non-English languages such as Korean \citep{koo2026emcee}.
To support this assessment, we follow the answer-conditioned rationalization approach of \citet{zelikman2022star} and use GPT-5.4~\citep{openai2026gpt54} to generate step-by-step solutions conditioned on the question, options, and gold answer, including explanations for incorrect options. 
We refer to these model-generated solutions as \textit{silver reasoning}, as they have not been verified by domain experts.

We assess answer recoverability on 150 sampled solutions (50 per benchmark; Appendix~\ref{app:silver}).
After removing explicit statements of the final answer from the silver reasoning, we ask a verifier (GPT-5.4) to answer each question given its options and silver reasoning. 
The verifier answers all 150 questions correctly, compared with 118 when given only the question and options. 
While this suggests that silver reasoning helps the verifier answer the questions correctly, it does not establish the factual correctness of individual reasoning steps.
We therefore use silver reasoning as supplementary reference material for the annotations below and as a source for deriving anchor knowledge in Sec.~\ref{sec:klear}, rather than as ground-truth solutions. 

\subsection{Knowledge and Application Demands}
\label{sec:ka}
\begin{table}[t]
    \centering
    \small
    \caption{\textbf{Validation of knowledge (K) and application (A) labels.} LLM--human agreement on 60 questions, averaged across two human annotators: exact match (EM; \%) and linearly weighted Cohen's $\kappa$, which penalizes disagreements in proportion to the distance between levels. \textit{w/o silver} and \textit{w/ silver} denote the LLM annotation without and with silver reasoning.}
    \label{tab:ka_val}

    \begin{tabular}{l|cc|cc}
        \toprule[1.5pt]
        & \multicolumn{2}{c|}{\textbf{K}} & \multicolumn{2}{c}{\textbf{A}} \\
        \cmidrule(lr){2-3} \cmidrule(lr){4-5}
        & EM & $\kappa$ & EM & $\kappa$ \\
        \midrule[1pt]
        w/o silver & 68.3 & 0.67 & 45.0 & 0.47 \\
        \rowcolor[HTML]{D9D9D9}
        w/ silver & 56.7 & 0.56 & 69.2 & 0.69 \\
        \bottomrule[1.5pt]
    \end{tabular}
    \vspace{-0.1in}
\end{table}

We characterize each question along two axes.
\begin{itemize}[leftmargin=3.5mm, topsep=2pt, itemsep=2pt, parsep=0pt]
    \item[$\circ$] \textit{Knowledge demand} (K) is the level of prior knowledge needed to solve the question: none (K0), everyday common sense (K1), school-level knowledge (K2), or specialized professional knowledge (K3).
    \item[$\circ$] \textit{Application demand} (A) is the complexity of connecting that knowledge to the case described in the question. 
    At A0, the retrieved fact or rule is itself the answer. 
    At A1, the correspondence between the knowledge and the case is stated explicitly.
    At A2, a single piece of knowledge applies, but one must judge whether the case satisfies its condition.
    At A3, two or more candidate rules compete, and one must decide which governs the case.
\end{itemize}
Both axes describe the question and its solution, independent of how the model performs.
We use GPT-5.4 to assign K and A, given the question, options, gold answer, and silver reasoning.
See Tables~\ref{tab:k_examples} and~\ref{tab:a_examples} for examples.

To validate these labels, two human annotators independently labeled 60 questions sampled evenly across the four A levels, seeing only the question, options, and gold answer (Appendix~\ref{app:ka_annotation}).
Table~\ref{tab:ka_val} compares their labels with those of the model, with and without silver reasoning.
With silver reasoning, the model agrees with the annotators on A ($\kappa=0.69$) more closely than the annotators agree with each other ($\kappa=0.60$).
Without silver reasoning, agreement on A drops to 0.47. 
Silver reasoning does not improve agreement on K, where the model tends to assign K2 to questions that annotators consider K3. 
Since application demand is the axis our analysis and construction focus on, we adopt silver reasoning for annotation.

\begin{table}[t]
    \centering
    \caption{\textbf{Distribution of knowledge (K) and application (A) demands.} Percentage of the 7,340 questions across K and A levels in three Korean benchmarks.}
    \label{tab:ka_dist}

    \resizebox{0.8\columnwidth}{!}{%
        \begin{tabular}{l|cccc|c}
            \toprule[1.5pt]
            & \textbf{A0} & \textbf{A1} & \textbf{A2} & \textbf{A3} & \textbf{Total} \\
            \midrule[1pt]
            \textbf{K0} & 0.0  & 1.0  & 0.3  & 0.0  & 1.3  \\
            \textbf{K1} & 2.9  & 1.2  & 1.1  & 0.1  & 5.2  \\
            \textbf{K2} & 18.5 & 4.0  & 5.5  & 3.6  & 31.7 \\
            \textbf{K3} & 39.6 & 4.2  & 5.1  & 12.7 & 61.7 \\
            \midrule
            \rowcolor[HTML]{D9D9D9}
            \textbf{Total} & 61.1 & 10.4 & 12.0 & 16.4 & 100.0 \\
            \bottomrule[1.5pt]
        \end{tabular}%
    }
    \vspace{-0.1in}
\end{table}

Table~\ref{tab:ka_dist} shows the resulting distribution.
Most questions demand specialized knowledge but little application:
61.7\% are K3 and 61.1\% are A0, and K3 questions answerable by recall alone account for 39.6\% of the total.
A0 is the largest level in every benchmark, ranging from 50.6\% in KMMLU-Pro to 73.6\% in KMMLU-Redux.
Questions that require application are fewer and unevenly distributed, with 72.8\% of A3 questions coming from KMMLU-Pro.

\begin{figure*}[t]
\centering
\includegraphics[width=\textwidth]{Figures/raw_figures/figure1.pdf}
%\vspace{-0.2in}
\caption{
\textbf{Overview of \name{}.}
\textbf{(a)} Seed questions from Korean benchmarks are labeled by knowledge (K) and application (A) demands.
\textbf{(b)} Silver reasoning is generated from the seed question and gold answer, then used to derive an anchor knowledge bundle for generating higher-A questions.
\textbf{(c)} Generated questions undergo LLM auditing, proxy screening, and human validation, with failed candidates revised using LLM or human feedback.
The result is a matched set of questions that share the same knowledge but differ in application demand.
}\label{fig:overview}
\vspace{-0.1in}
\end{figure*}

\subsection{Failure Analysis}
\label{sec:failure}

We next examine how models fail on these questions.
We evaluate three Korean LLMs, K-EXAONE \citep{choi2026k}, A.X K1 \citep{cheon2026ax}, and Solar Open \citep{park2026solar}, with reasoning enabled, and collect their 5,808 incorrect responses.
For each response, GPT-5.4 reads the model's reasoning trace alongside the silver reasoning and assigns one or more failure types, ordered by their contribution to the error (Appendix~\ref{app:source_failure}).
We group these types into two categories.
\begin{itemize}[leftmargin=3.5mm, topsep=2pt, itemsep=2pt, parsep=0pt]
    \item[$\circ$] \textit{Knowledge failure}: the model recalls an incorrect fact, attributes a fact to the wrong entity, or never engages the required rule.
    \item[$\circ$] \textit{Application failure}: the model invokes the relevant knowledge but misapplies it, for example by misjudging a condition or an exception.
\end{itemize}
The remaining responses stem from invalid outputs or defective questions.
We classify each response by its primary failure type. 

\begin{figure}[t]
\centering
\includegraphics[width=\columnwidth]{Figures/raw_figures/figure2.pdf}
\vspace{-0.2in}
\caption{
\textbf{Failure composition by application demand.}
Distribution of incorrect responses by failure type across three Korean LLMs, with accuracy (\%) on the right.
}\label{fig:failure}
\vspace{-0.15in}
\end{figure}

Figure~\ref{fig:failure} shows that failure composition differs between questions answerable by recall alone and questions that require application.
On A0 questions, knowledge failures dominate (69.8\% vs.\ 21.1\% application failures).
On A1--A3 questions, application failures account for 36.6--44.0\% of incorrect responses, comparable to knowledge failures (40.5--44.8\%).
The same pattern holds among K3 questions alone (Appendix~\ref{app:source_failure}).
Among questions that require application, accuracy also falls from 88.5\% at A1 to 60.0\% at A3.

These comparisons, however, are not controlled, because questions at different A levels also differ in the knowledge they require.
For instance, the share of K3 questions ranges from 40.3\% at A1 to 77.3\% at A3, and a single benchmark of professional licensure exams accounts for 72.8\% of A3 questions but at most 33.4\% at any other level. 
A gap between A levels may thus reflect what the questions are about rather than how much application they demand.
Isolating the effect of application demand requires contrasts in which the required knowledge is held fixed, which we construct in Sec.~\ref{sec:klear}.